% ============================ SOVEREIGN PRESS - shared preamble ============================
% One preamble for every book in the press. Compile with xelatex (fonts are set
% by fontspec via the pandoc template). Pandoc inlines this file into each book's
% generated .tex with -H, so its path does not matter at LaTeX time. Two things
% DO need to be found on the TeX input path at compile time, and build.py puts
% them there via TEXINPUTS: "version.tex" (the generated per-build tag/header
% snippet) and the book's images (from the book's assets/ and shared/assets/).
%
% This preamble is the series styling. It is deliberately book-agnostic: anything
% that differs per book (the version tag, the running-header label) arrives as a
% macro that build.py writes into version.tex. The workbook tooling at the bottom
% (checkboxes, the three-track command box, prompt panels, do-it lists) is what the
% Primer uses to become a workbook; the Atlas never calls it, so it changes nothing there.

% ---- Page geometry: A4 ----
\usepackage{geometry}
\def\sspaperclass{a4}
\geometry{a4paper,top=2.6cm,bottom=2.5cm,left=2.7cm,right=2.7cm,headheight=15pt}
\def\sstagx{20cm}\def\sstagy{28.7cm}

% ---- Unicode glyph fallbacks (defensive; manuscripts are pure ASCII) ----
\usepackage{newunicodechar}
\newunicodechar{→}{\ensuremath{\rightarrow}}
\newunicodechar{≈}{\ensuremath{\approx}}
\newunicodechar{…}{\dots}

% ---- Colours (deep green: the series accent) ----
\usepackage{xcolor}
\definecolor{accent}{HTML}{1B5E3A}
\definecolor{accentlt}{HTML}{2E7D4F}
\definecolor{codebg}{HTML}{F6F6F3}
\definecolor{codeframe}{HTML}{D7D7CE}
\definecolor{codecmnt}{HTML}{6A7A6A}
\definecolor{rulegray}{HTML}{C2C2BA}
\definecolor{softblack}{HTML}{1A1A1A}
% Track accents for the three-track command box (linux/mac/win), plus trkweb, the
% prompt-panel accent (not a track); muted, distinct, print-safe.
\definecolor{trklinux}{HTML}{1B5E3A}
\definecolor{trkmac}{HTML}{4A4A4A}
\definecolor{trkwin}{HTML}{1F5673}
\definecolor{trkweb}{HTML}{8A5A1B}

% ---- graphics (book screenshots, cover wash) ----
\usepackage{graphicx}
\graphicspath{{./}}            % real image dirs are supplied via TEXINPUTS by build.py
\usepackage{eso-pic}           % single-page title background

% ---- version tag + running-header label, generated per build ----
% build.py writes temp/<slug>/version.tex with \ssversion and \ssrunhead, and
% adds that folder to TEXINPUTS, so the tag and header update themselves on every
% build with no manual edit. The placeholders below are used only if a .tex is
% compiled directly (without build.py); i00 flags "version unknown".
\def\ssversion{v0v0i00}
\def\ssrunhead{\color{accent}\textbf{SOVEREIGN STACK}}
\IfFileExists{version.tex}{%% generated by build.sh (src/mkversion.py); do not edit by hand
\def\ssversion{v1v1i32}
}{}
\AddToShipoutPictureFG{\put(\LenToUnit{\sstagx},\LenToUnit{\sstagy}){\makebox(0,0)[r]{\ttfamily\footnotesize\color{softblack!50}\ssversion}}}

% ---- TikZ (the Atlas roadmap diagrams; the workbook checkbox) ----
\usepackage{tikz}
\usetikzlibrary{positioning,fit,backgrounds,arrows.meta}

% ---- Heading design (titlesec) ----
\usepackage{titlesec}
\titleformat{\chapter}[display]
  {\sffamily\bfseries\color{accent}}{}{0pt}{\Huge\filright}
  [\vspace{5pt}{\color{rulegray}\titlerule[1.2pt]}]
\titlespacing*{\chapter}{0pt}{6pt}{26pt}
\titleformat{\section}
  {\sffamily\bfseries\large\color{accentlt}}{}{0pt}{}
\titlespacing*{\section}{0pt}{18pt}{5pt}
\titleformat{\subsection}
  {\sffamily\bfseries\normalsize\color{softblack}}{}{0pt}{}
\titlespacing*{\subsection}{0pt}{12pt}{3pt}

% ---- Running headers / footers ----
\usepackage{fancyhdr}
\usepackage{lastpage}          % \pageref{LastPage}: total page count for the footer
\pagestyle{fancy}
\fancyhf{}
\renewcommand{\chaptermark}[1]{\markboth{#1}{}}
\fancyhead[L]{\footnotesize\sffamily\ssrunhead}
\fancyhead[R]{\footnotesize\sffamily\color{softblack}\nouppercase{\leftmark}}
\fancyfoot[C]{\small\sffamily\thepage\ /\ \pageref{LastPage}}
\renewcommand{\headrulewidth}{0.4pt}
\renewcommand{\footrulewidth}{0pt}
\fancypagestyle{plain}{\fancyhf{}\fancyfoot[C]{\small\sffamily\thepage\ /\ \pageref{LastPage}}%
  \renewcommand{\headrulewidth}{0pt}}

% ---- Code listings (wrap long lines so nothing overflows A4) ----
\usepackage{listings}
\lstset{
  basicstyle=\ttfamily\small,
  commentstyle=\color{codecmnt}\itshape,
  backgroundcolor=\color{codebg},
  frame=single,
  rulecolor=\color{codeframe},
  framesep=7pt,
  framexleftmargin=2pt,
  xleftmargin=12pt,
  xrightmargin=6pt,
  breaklines=true,
  breakatwhitespace=false,
  prebreak=\mbox{\textcolor{codeframe}{\tiny$\searrow$}},
  postbreak=\mbox{\textcolor{codeframe}{\small$\hookrightarrow$}\space},
  keepspaces=true,
  showstringspaces=false,
  extendedchars=true,
  columns=fullflexible,
  aboveskip=12pt,
  belowskip=12pt
}

% ---- Wide data tables (e.g. the Atlas bandwidth timeline): shrink + tighten ----
\usepackage{etoolbox}
\usepackage{array}
\AtBeginEnvironment{longtable}{\footnotesize\setlength{\tabcolsep}{4.5pt}}
\renewcommand{\arraystretch}{1.12}

% ---- Figure-fit helper (wraps the Atlas's wide TikZ diagrams) ----
% A plain pass-through at A4 (the wide boxes already fit the text width). Kept so
% the manuscript's \ssfit{...} calls keep compiling; if a narrower edition is ever
% added, the scaling logic returns here without editing any book.
\newcommand{\ssfit}[1]{#1}

% ---- Pointer / callout boxes (markdown blockquotes -> accent-barred callouts) ----
\usepackage[most]{tcolorbox}
\renewenvironment{quote}{%
  \begin{tcolorbox}[breakable, enhanced, colback=codebg, colframe=accentlt,
    boxrule=0pt, leftrule=2.6pt, arc=0pt, outer arc=0pt,
    left=11pt, right=11pt, top=8pt, bottom=8pt,
    before skip=11pt, after skip=11pt]\small\color{softblack}}{%
  \end{tcolorbox}}

% ---- Numbering / TOC ----
% Headings carry their own labels ("PART 0", "0.1 ..."), so suppress LaTeX's
% automatic section numbers but keep everything in the TOC.
\setcounter{secnumdepth}{-1}
\setcounter{tocdepth}{1}

% ---- Spacing, lists, URL breaking ----
\usepackage{microtype}
\usepackage{xurl}
\usepackage{enumitem}
\setlist{itemsep=2pt,topsep=4pt,parsep=0pt}
\setlength{\parindent}{0pt}
\setlength{\parskip}{0.55em plus 0.1em}
\sloppy
\emergencystretch=2.5em

% ============================ WORKBOOK TOOLING ============================
% Everything below exists for the Primer (the slow, do-it workbook). The Atlas
% never calls any of it. It gives four things: a pen-tickable checkbox, a
% three-track command box (one row per kind of computer, reserved for exact
% platform-specific instructions), a prompt panel (natural language destined
% for a search box or an AI chat, the same on every machine), and a "Do it"
% task list whose every line begins with a checkbox. Together they turn a
% concept from something you read into one page you can work through and tick.

% A small, pen-tickable square in the accent colour.
\newcommand{\ssbox}{\tikz[baseline=-0.45ex]\draw[accent,line width=0.7pt,rounded corners=0.6pt] (0,0) rectangle (1.05ex,1.05ex);}

% One row of a three-track box: a fixed-width coloured track label, then the
% command (verbatim-ish) or a note. Use \track for a literal command, \tracknote
% for prose (e.g. the phone track, where the answer is not a command).
\newcommand{\tracklabelbox}[2]{\makebox[2.9cm][l]{\small\sffamily\bfseries\color{#1}#2}}
\newcommand{\track}[3]{\tracklabelbox{#1}{#2}\texttt{\small #3}\par}
\newcommand{\tracknote}[3]{\tracklabelbox{#1}{#2}{\small\itshape\color{softblack!80}#3}\par}
% Convenience wrappers naming the three canonical tracks, so a book never has
% to remember the colour. CachyOS/Linux, macOS, Windows.
\newcommand{\tlinux}[1]{\track{trklinux}{CachyOS}{#1}}
\newcommand{\tmac}[1]{\track{trkmac}{macOS}{#1}}
\newcommand{\twin}[1]{\track{trkwin}{Windows}{#1}}
\newcommand{\tlinuxn}[1]{\tracknote{trklinux}{CachyOS}{#1}}
\newcommand{\tmacn}[1]{\tracknote{trkmac}{macOS}{#1}}
\newcommand{\twinn}[1]{\tracknote{trkwin}{Windows}{#1}}

% The three-track box itself: a titled, tinted panel that holds the rows above.
\newtcolorbox{tracks}{enhanced, breakable, colback=codebg, colframe=rulegray,
  boxrule=0.6pt, arc=1pt, left=10pt, right=10pt, top=7pt, bottom=7pt,
  before skip=10pt, after skip=10pt,
  fonttitle=\sffamily\bfseries\small\color{softblack},
  coltitle=softblack, title={On your machine}}

% The prompt panel: natural language to type into a search box or an AI chat.
% Deliberately unlike the tracks box: no per-machine rows, because the words
% are the same everywhere; slightly tinted, body in italic to signal "this is
% language, not a command". CAPITALISED words inside are placeholders the
% reader swaps for their own details.
\newtcolorbox{promptbox}{enhanced, breakable, colback=trkweb!6!white,
  colframe=trkweb!70!rulegray, boxrule=0.6pt, arc=1pt,
  left=10pt, right=10pt, top=7pt, bottom=7pt,
  before skip=10pt, after skip=10pt,
  fonttitle=\sffamily\bfseries\small, coltitle=softblack,
  fontupper=\small\itshape, title={The prompt}}

% The "Do it" task list: a tinted panel whose every item starts with a checkbox.
% Two ways to fill it. (1) The doitlist environment, for hand-written LaTeX. (2)
% A plain itemize, for Markdown-authored pages: the box redefines the first-level
% itemize label to the checkbox, so a normal Markdown bullet list written inside
% a "::: doit" block (which pandoc turns into an itemize) comes out ticked too.
\newlist{doitlist}{itemize}{1}
\setlist[doitlist]{label=\ssbox, leftmargin=1.7em, itemsep=4pt, topsep=2pt}
\newtcolorbox{doit}{enhanced, breakable, colback=white, colframe=accent,
  boxrule=0.8pt, arc=1pt, left=10pt, right=10pt, top=7pt, bottom=7pt,
  before skip=10pt, after skip=10pt,
  before upper={\renewcommand{\labelitemi}{\ssbox}\setlength{\leftmargini}{1.7em}},
  fonttitle=\sffamily\bfseries\small\color{white},
  coltitle=white, title={Do it}}

% The success line: what the reader should see or feel when the page worked.
\newcommand{\seeline}[1]{\par\vspace{3pt}{\small\sffamily\color{accentlt}\textbf{You should see\;}}{\small\color{softblack}#1}\par}

% A workbook concept heading: starts a fresh page, records a Contents entry, and
% prints the number and title over a rule. This is the per-page header the
% Markdown "### N. Title" lines become (one concept to a page). The explanation,
% the three-track box, the do-it list and the success line then follow as ordinary
% content, so the page is assembled from plain Markdown plus a few small blocks.
\newcommand{\concept}[2]{%
  \clearpage
  \phantomsection\addcontentsline{toc}{section}{#1\hspace{0.6em}\textbar\hspace{0.6em}#2}%
  {\sffamily\bfseries\Large\color{accent}#1\quad\color{softblack}#2\par}
  \vspace{2pt}{\color{rulegray}\hrule height 0.8pt}\vspace{8pt}}

% A single workbook page assembled in one call, for hand-written LaTeX. Arguments:
%   {number}{title}{image path or empty}{plain explanation}
%   {track rows}{do-it items}{success line}
% Any argument may be left empty. The Markdown "::: " block syntax (see
% preprocess.py) is the usual way to author pages; this macro is the all-in-one
% alternative when writing LaTeX directly.
\newcommand{\workpage}[7]{%
  \clearpage
  \phantomsection\addcontentsline{toc}{section}{#1\hspace{0.6em}\textbar\hspace{0.6em}#2}%
  {\sffamily\bfseries\Large\color{accent}#1\quad\color{softblack}#2\par}
  \vspace{2pt}{\color{rulegray}\hrule height 0.8pt}\vspace{8pt}
  \ifstrempty{#3}{}{\begin{center}\includegraphics[width=0.62\textwidth]{#3}\end{center}\vspace{4pt}}%
  #4\par
  \ifstrempty{#5}{}{\begin{tracks}#5\end{tracks}}%
  \ifstrempty{#6}{}{\begin{doit}\begin{doitlist}#6\end{doitlist}\end{doit}}%
  \ifstrempty{#7}{}{\seeline{#7}}%
  \clearpage}
